У ході роботи розроблено метод автоматичної ідентифікації музичних інструментів в аудіозаписах на основі гібридної архітектури CNN~+~BiLSTM з мел-спектрограмами у постановці multi-label класифікації.

Основні результати роботи:
\begin{itemize}
  \item проведено огляд і порівняльний аналіз існуючих методів автоматичної ідентифікації музичних інструментів; обґрунтовано перевагу гібридних CNN~+~BiLSTM-архітектур над класичними ML-методами та виключно CNN-підходами;
  \item побудовано математичну модель методу, що описує повний цикл перетворення аудіосигналу у вектор інструментальних міток: від мел-спектрограми через CNN і BiLSTM до незалежних сигмоїдних активацій та бінарної крос-ентропійної функції втрат;
  \item реалізовано програмне забезпечення методу мовою Python~3 з використанням PyTorch і librosa;
  \item метод верифіковано та валідовано на датасеті IRMAS (11 класів інструментів, 6000 фрагментів); отримані значення F1-macro, F1-micro та mAP перевищують результати базових методів.
\end{itemize}

Отримані результати підтверджують доцільність застосування мел-спектрограм як тембрального представлення та гібридної CNN~+~BiLSTM архітектури для задачі multi-label ідентифікації інструментального складу аудіозаписів. Розроблений метод може бути використаний для побудови систем автоматичної анотації та контентного пошуку у музичних колекціях.
